\section*{अध्याय \ref{ch:g1:materials-around-us} --- चीज़ें किससे बनी हैं}

\begin{solution}{exo:g1:materials-around-us:1}
कुर्सी एक \omterm{def:g1:materials-around-us:object}{वस्तु} है। लकड़ी एक \omterm{def:g1:materials-around-us:material}{सामग्री} है: उससे कुर्सी बन सकती है।
\end{solution}

\begin{solution}{exo:g1:materials-around-us:2}
खिड़की का शीशा: काँच। कील: धातु। कंकड़: पत्थर। किताब: कागज़। मोज़ा:
कपड़ा (ऊन या सूत)।
\end{solution}

\begin{solution}{exo:g1:materials-around-us:3}
लकड़ी: पेंसिल, दरवाज़ा (या कुर्सी, मेज़)। प्लास्टिक: कटोरी, बाल्टी
(या कोई खिलौना, बोतल)।
\end{solution}

\begin{solution}{exo:g1:materials-around-us:4}
धातु का काँटा: बस वही लकड़ी या प्लास्टिक का नहीं बना है। (दो लकड़ी
के बने हैं और एक प्लास्टिक का।)
\end{solution}

\begin{solution}{exo:g1:materials-around-us:5}
काँच: वह पारदर्शी है।
\end{solution}

\begin{solution}{exo:g1:materials-around-us:6}
स्टील की कील। चुंबक लोहे और स्टील को खींचता है, काँच या लकड़ी को नहीं।
\end{solution}

\begin{solution}{exo:g1:materials-around-us:7}
दो वस्तुएँ तैरती हैं (लकड़ी का गुटका और प्लास्टिक का ढक्कन); तीन डूबती
हैं (कंकड़, कंचा और कील)।
\end{solution}

\begin{solution}{exo:g1:materials-around-us:8}
फल धातु के हैं, हत्थे प्लास्टिक के। कैंची दो सामग्रियों से
बनी है।
\end{solution}

\begin{solution}{exo:g1:materials-around-us:9}
नीले घेरे में: वह धातु का बना है। वह नारंगी घेरे में नहीं है,
क्योंकि चुंबक उसे नहीं खींचता: वह ऐलुमिनियम का बना है, लोहे या
स्टील का नहीं।
\end{solution}

\begin{solution}{exo:g1:materials-around-us:10}
खिड़की को प्रकाश अंदर आने देना है, और हमें उसके आर-पार देखना है। काँच
पारदर्शी है; लकड़ी नहीं।
\end{solution}

\begin{solution}{exo:g1:materials-around-us:11}
``हाँ'': कील, पेपर-क्लिप और चाबी, 3 वस्तुएँ। ``नहीं'': कंकड़,
चम्मच, डिब्बा, कंचा और मोज़ा, 5 वस्तुएँ।
$3 + 5 = 8$: हर \omterm{def:g1:materials-around-us:object}{वस्तु} को उसकी जगह मिल गई।
\end{solution}
